\documentclass[10pt,a4paper]{amsart}
\usepackage[margin=19mm]{geometry}
\usepackage{amsmath,amssymb,mathtools}
\usepackage[scr=rsfs,cal=euler]{mathalfa}
\usepackage{xfrac}
\usepackage[unicode,hidelinks,bookmarks=false]{hyperref}
\newcommand{\ie}{i.e.\ }
\newcommand{\cf}{cf.\ }
\newcommand{\resp}{respectively\ }
\newcommand{\Subsection}{\subsection}
\providecommand{\nobreakdash}{\nobreak\hbox{-}}
\setlength{\emergencystretch}{2em}
\multlinegap0pt
\newcommand{\Disc}{\operatorname{Disc}}
\newcommand{\Res}{\operatorname{Res}}
\theoremstyle{plain}
\newtheorem{theorem}{Theorem}[section]
\newtheorem{proposition}[theorem]{Proposition}
\newtheorem{lemma}[theorem]{Lemma}
\newtheorem{corollary}[theorem]{Corollary}
\theoremstyle{definition}
\newtheorem{remark}[theorem]{Remark}
\newtheorem{example}[theorem]{Example}
\title[Resultant of an equivariant polynomial system with respect to the reflection group $G(r,n)$]{Resultant of an equivariant polynomial system with respect to the reflection group $G(r,n)$}
\author{Sonagnon Julien Owolabi, Ibrahim Nonkan\'{e}, Joel Tossa}
\date{Source: arXiv:2609.27886v1 (CC BY 4.0)}
\begin{document}
\maketitle

\noindent\emph{Source and licence.} The delivered work is the article shown in the title above, version arXiv:2609.27886v1 (CC BY 4.0), distributed under the Creative Commons Attribution 4.0 International licence, as recorded in the arXiv licence metadata for this exact version and shown on the abstract page saved with this item. The full licence notice, which carries the licence URL and the address of that abstract page, is the bundled file \texttt{licenses/arxiv-2609.27886-CC-BY-4.0.txt}.
\par\smallskip

\noindent\textit{Editorial scope.} Counted unit: the article's Theorem (source label thm:disc2) (source0.tex lines 573--580, source label \texttt{thm:disc2}): ``Let of degree be a homogeneous polynomial invariant under . Assume , , , and let , , be as in Proposition . Then d a(n,d) (f) r n(d-1) n-1 ( f 1 , , f n ) i 1 n-1 ( f 1 , , f i , x i 1 , ,x n ) n i (r''. It is delivered together with the article's complete original proof (source0.tex lines 582--615, one proof environment adjacent to the statement, 34 lines), copied line for line from the article's own text. Required context retained uncounted: lem:powersubstitution (lines 224--230); eq:discres (lines 491--494); prop:theta (lines 555--562). Editorial formatting only: an AMS \texttt{amsart} preamble that restates the article's own macro definitions and its own theorem declarations, and the theorem environments the retained lines use; the article's own class file is neither shipped nor input; the retained environments are renumbered sequentially in order of appearance, whereas the article prints them with its own numbering; the article's equation numbering is not reproduced and every cross-reference inside the retained text resolves to the wrapper's numbering. No citation occurs inside any retained line, so the wrapper carries no bibliography; All retained bodies are exact copies of the bundled \texttt{sources/211522/source0.tex}, so no omission receipt applies. No asset-inclusion command occurs inside any retained span and the package ships no image asset.


\section*{Retained context: lem:powersubstitution (source0.tex lines 224--230)}

\begin{lemma}
\label{lem:powersubstitution}
Let $n \geqslant 2$ and $r \geqslant 1$ be integers, and let $Q_1,\dots,Q_n \in \mathbb{C}[y_1,\dots,y_n]$ be homogeneous polynomials of respective degrees $d_1',\dots,d_n' \geqslant 1$, not necessarily equal. Set $f_i(x_1,\dots,x_n) := Q_i(x_1^r,\dots,x_n^r)$, homogeneous of degree $r d_i'$. Then
\[
\Res(f_1,\dots,f_n) = \Res(Q_1,\dots,Q_n)^{r^{n-1}}.
\]
\end{lemma}


\section*{Retained context: eq:discres (source0.tex lines 491--494)}

\begin{equation}
\label{eq:discres}
d^{a(n,d)}\Disc(f) = \Res\big(f^{\{1\}},f^{\{2\}},\dots,f^{\{n\}}\big),
\end{equation}


\section*{Retained context: prop:theta (source0.tex lines 555--562)}

\begin{proposition}
\label{prop:theta}
Let $f \in \mathbb{C}[x_1,\dots,x_n]$ of degree $d$ be a homogeneous polynomial invariant under the reflection group $G(r,n)$. For $r>1$ and for all $i=1,\dots,n$, there exists a homogeneous polynomial $\Theta f^{\{i\}}$ of degree $d-r$ such that
\[
\frac{\partial f}{\partial x_i}(x_1,\dots,x_n) = r x_i^{r-1}\, \Theta f^{\{i\}}(x_1,\dots,x_n).
\]
The polynomials $\Theta f^{\{1\}},\Theta f^{\{2\}},\dots,\Theta f^{\{n\}}$ form a $G(r,n)$-equivariant system of homogeneous polynomials of degree $d-r$.
\end{proposition}


\section{The counted unit: Theorem (source label thm:disc2) and its complete original proof}

\begin{theorem}
\label{thm:disc2}
Let $f \in \mathbb{C}[x_1,\dots,x_n]$ of degree $d$ be a homogeneous polynomial invariant under $G(r,n)$. Assume $n\geqslant 2$, $d\geqslant 2$, $r>1$, and let $\Theta f^{\{i\}}$, $i=1,\dots,n$, be as in Proposition~\ref{prop:theta}. Then
\begin{multline*}
d^{a(n,d)}\Disc(f) = r^{n(d-1)^{n-1}} \times \Res\big(\Theta f^{\{1\}},\dots,\Theta f^{\{n\}}\big) \\
\times \prod_{i=1}^{n-1} \Res\big(\Theta f^{\{1\}},\dots,\Theta f^{\{i\}}, x_{i+1},\dots,x_n\big)^{\binom{n}{i}(r-1)^{n-i}}.
\end{multline*}
\end{theorem}

\begin{proof}
By \eqref{eq:discres} and Proposition~\ref{prop:theta},
\begin{align*}
d^{a(n,d)}\Disc(f) &= \Res\left(\frac{\partial f}{\partial x_1},\dots,\frac{\partial f}{\partial x_n}\right) = \Res\big(rx_1^{r-1}\Theta f^{\{1\}},\dots,rx_n^{r-1}\Theta f^{\{n\}}\big) \\
&= r^{n(d-1)^{n-1}}\Res\big(x_1^{r-1}\Theta f^{\{1\}},\dots,x_n^{r-1}\Theta f^{\{n\}}\big).
\end{align*}
By multiplicativity of the resultant, applied once in each of the $n$ slots, $\Res\big(x_1^{r-1}\Theta f^{\{1\}},\dots,x_n^{r-1}\Theta f^{\{n\}}\big)$ splits into a product of $2^n$ factors $\Res(h_1,\dots,h_n)$, indexed by the subsets $S \subseteq \{1,\dots,n\}$ of slots where $h_i := x_i^{r-1}$ is chosen (and $h_i := \Theta f^{\{i\}}$ for $i \notin S$).

These $2^n$ factors depend only on $|S|$. Fix $i := n-|S|$, the number of slots carrying a $\Theta f$, and set $D := (r-1)^{|S|}(d-r)^{n-|S|}$, the product of the degrees of $h_1,\dots,h_n$ (depending only on $|S|$, not on which slots are in $S$). Given two subsets $S,S'\subseteq\{1,\dots,n\}$ of the same size, choose $\sigma \in S_n$ with $\sigma(\{1,\dots,n\}\setminus S) = \{1,\dots,n\}\setminus S'$. Since $\Theta f^{\{1\}},\dots,\Theta f^{\{n\}}$ form a $G(r,n)$-equivariant system (Proposition~\ref{prop:theta}), they are in particular equivariant under the pure permutation part of $G(r,n)$, i.e.\ $\Theta f^{\{k\}}(x_{\sigma(1)},\dots,x_{\sigma(n)}) = \Theta f^{\{\sigma(k)\}}(x_1,\dots,x_n)$ for every $k$. Substituting $x_k \mapsto x_{\sigma(k)}$ in every entry of the tuple $(h_1,\dots,h_n)$ indexed by $S$ sends slot $i$ to $x_{\sigma(i)}^{r-1}$ (if $i\in S$) or to $\Theta f^{\{\sigma(i)\}}$ (if $i\notin S$); relabelling the slots by $j=\sigma(i)$ shows that, up to reordering the $n$ arguments by $\sigma$, this is exactly the tuple indexed by $\sigma(S)$, which equals $S'$. By the ``linear change of variables'' property, the substitution $x_k \mapsto x_{\sigma(k)}$ multiplies the resultant by $\varepsilon(\sigma)^D$; by the ``permutation of polynomials'' property, reordering the $n$ arguments by $\sigma$ multiplies it by the same factor $\varepsilon(\sigma)^D$ again. The two signs multiply to $\varepsilon(\sigma)^{2D}=1$, so
\[
\Res(h_1,\dots,h_n)\big|_{\text{indexed by }S} = \Res(h_1,\dots,h_n)\big|_{\text{indexed by }S'}.
\]
Consequently all $\binom{n}{i}$ subsets $S$ of size $n-i$ give the same value, which we compute using the canonical choice $S=\{i+1,\dots,n\}$:
\[
R_i := \Res\big(\Theta f^{\{1\}},\dots,\Theta f^{\{i\}}, x_{i+1}^{r-1},\dots,x_n^{r-1}\big), \qquad i=0,1,\dots,n,
\]
(with the convention that no $\Theta f$ appears when $i=0$, and no $x^{r-1}$ appears when $i=n$), so that
\[
\Res\big(x_1^{r-1}\Theta f^{\{1\}},\dots,x_n^{r-1}\Theta f^{\{n\}}\big) = \prod_{i=0}^n R_i^{\binom{n}{i}}.
\]
We have $R_0 = \Res(x_1^{r-1},\dots,x_n^{r-1}) = 1$ by the normalization of the resultant, and $R_n = \Res\big(\Theta f^{\{1\}},\dots,\Theta f^{\{n\}}\big)$.

\emph{Reduction of $R_i$ for $1\leqslant i \leqslant n-1$.} The map $(x_1,\dots,x_n) \mapsto (x_1,\dots,x_i,x_{i+1}^{r-1},\dots,x_n^{r-1})$ is the identity on the first $i$ coordinates and the power map $x\mapsto x^{r-1}$ on the last $n-i$; it defines a morphism $\mathbb{P}^{n-1}(\mathbb{C})\to\mathbb{P}^{n-1}(\mathbb{C})$ onto of degree $(r-1)^{n-i}$. Exactly the same argument as in the proof of Lemma~\ref{lem:powersubstitution} (same vanishing locus, proportionality by irreducibility of the universal resultant, degree count via the classical homogeneity property, and normalization by specializing $\Theta f^{\{k\}}=y_k$) gives
\begin{align*}
R_i &= \Res\big(\Theta f^{\{1\}},\dots,\Theta f^{\{i\}},x_{i+1}^{r-1},\dots,x_n^{r-1}\big) \\
&= \Res\big(\Theta f^{\{1\}},\dots,\Theta f^{\{i\}},x_{i+1},\dots,x_n\big)^{(r-1)^{n-i}}.
\end{align*}
Altogether,
\begin{align*}
\Res\big(x_1^{r-1}\Theta f^{\{1\}},\dots,x_n^{r-1}\Theta f^{\{n\}}\big) &= \prod_{i=1}^{n-1}\Res\big(\Theta f^{\{1\}},\dots,\Theta f^{\{i\}},x_{i+1},\dots,x_n\big)^{\binom{n}{i}(r-1)^{n-i}} \\
&\quad\times \Res\big(\Theta f^{\{1\}},\dots,\Theta f^{\{n\}}\big),
\end{align*}
and we deduce the claimed formula.
\end{proof}

\end{document}
